\section{微分すると簡単になる因子を選ぶ}\label{節:部分積分の因子選択}

整式は微分すると次数が下がるため，微分する側の候補になる。指数関数や三角関数は，積分しても同じ種類の関数のままなので，積分する側に選びやすい。ただし，この順序は絶対の規則ではない。対数を含む積では，対数を微分すると分数になり，別の因子と約分できることがある。

また，見た目が積でなくても，$1$ を掛けた積と考えられる。$\log x$ の原始関数を直接思い出せない場合でも，$\log x\cdot1$ と見れば，積分する側に既知の関数 $1$ を置ける。

\begin{example}{一つの関数を積と見る}
$x>0$ において，$\int\log x\dd x$ を求めよ。
\end{example}

$\log x$ の微分は $\frac1x$ である。これに $1$ の原始関数 $x$ が掛かれば，積は $1$ になる。そこで
\[
 u=\log x,\qquad v'=1,
 \qquad u'=\frac1x,\qquad v=x
\]
と選ぶと，
\[
\begin{aligned}
 \int\log x\dd x
 &=x\log x-\int x\cdot\frac1x\dd x\\
 &=x\log x-x+C.
\end{aligned}
\]
もとの対数を，微分して得られる分数と，積分して得られる $x$ の組合せで消したのである。微分による確認は
\[
 \frac{\dd}{\dd x}(x\log x-x)
 =\log x+x\cdot\frac1x-1=\log x
\]
となる。$\log x$ は $x>0$ で考えているので，この計算にも同じ条件が付く。

この選び方は $x\log x$ にも使える。$\log x$ を微分し，$x$ を積分すると，
\[
\begin{aligned}
 \int x\log x\dd x
 &=\frac{x^2}{2}\log x
   -\int\frac{x^2}{2}\cdot\frac1x\dd x\\
 &=\frac{x^2}{2}\log x-\frac{x^2}{4}+C.
\end{aligned}
\]
ここでは整式 $x$ を積分する側に選んだ。それでも，対数を微分して得た $\frac1x$ と掛けることで，残る積分が単なる一次式になる。関数の名前だけで順序を決めず，$u'v$ を具体的に見ることが大切である。

一方，$\int\frac{\log x}{x}\dd x$ ならば，前章で見たとおり $t=\log x$ と置くほうが短い。$\frac1x\dd x=\dd t$ となり，ただちに $\int t\dd t$ へ変わるからである。「対数があるから部分積分」と決めるのではなく，対数の微分がすでに残りの部分にあるかを先に調べよう。

因子の選択には，もう一つ条件がある。積分する側の原始関数 $v$ を，その時点で求められなければならない。たとえば $\int x\e^{x^2}\dd x$ で，$x$ を微分しようとして $v'=\e^{x^2}$ と選ぶと，そこで原始関数が分からなくなる。しかし $t=x^2$ と置けば，$x\dd x=\frac12\dd t$ なので，答えは $\frac12\e^{x^2}+C$ である。部分積分を試す前に，置換に必要な微分の因子がそろっていないかを見ると，遠回りを避けられる。

